% The paid local components and chamber receipt assemble to a one-point probe.
\section{Assembly of the finite-volume response probe}
\label{sec:ex-lattice-trace-probe-assembly}

We spell out the finite algebra that turns two separately paid local
components and a buffered chamber approximation into the probe used in the
mission. The quantitative assumptions below are exactly the estimates the
assembly consumes. In particular, the proof establishes every component of
the probe conclusion and does not regard the existence of a record with
those components as a proof.

Fix a frame level $j$, its volume level $k$, and a finite volume $N>k$.
Write
\[
 t=\widehat\beta_k,\qquad a>0,\qquad \beta>0,\qquad
 q=(a/a_k)^5,
\]
where $a_k$ is the field scale at the volume level. In the selected
construction $k$ is the summit index of $j$ (so $k>j$), $a=a_j$ is the
field scale at the frame level, and $\beta$ is the finite-volume coupling
at level $k$ in the volume $N$. The field scales satisfy $a_0=1$ and
decrease strictly, so $a_j<1$ for $j>0$: only the positivity of $a$ is
available, and only positivity is used below. Let $\tau_j$ be the
error scale
\[
 \tau_j(s)=\exp\!\left(\frac{s}{10^4}
           (\log\widehat\beta_j)^{48}\right),
 \qquad E=\tau_j(-430).
\]
Assume the scale bounds
\begin{equation}
 t\geq1,\quad 0<\beta\leq2t,\quad
 1+2t\leq\tau_j(230),\quad 6c\leq t,
 \label{eq:ex-trace-probe-scales}
\end{equation}
where $c$ is the coefficient cost of the finite symmetry average defined
below. The error scales obey $\tau_j(s)\tau_j(s')=\tau_j(s+s')$.

Let $\mathfrak A=\operatorname{Signed}(4)\times\{0,1\}$ be the finite
motion set; the second factor records whether complex conjugation is
applied. For $\mu=(s,e)$, write $\rho_\mu$ for the real-linear action on
grade-five coefficients and $\sigma_\mu$ for the corresponding real-linear
action on observables. For $\mu=(s,e)$ and axis $d$, define
$\operatorname{prev}(\mu,d)=s.\operatorname{perm}^{-1}(d)$ and
$\epsilon(\mu,d)=\operatorname{sign}(s,\operatorname{prev}(\mu,d))$;
the sign does not depend on the conjugation bit $e$. Put
$R=|\mathfrak A|^{-1}$. The finite averages have the
form
\begin{align}
 (\mathscr S_5 V)_d&=R\sum_{\mu\in\mathfrak A}
     \epsilon(\mu,d)\,\rho_\mu V_{\operatorname{prev}(\mu,d)},
 &R\sum_{\mu\in\mathfrak A}|\epsilon(\mu,d)|&=1,
 \label{eq:ex-trace-probe-coeff-average}\\
 (\mathscr S_0 f)_d&=R\sum_{\mu\in\mathfrak A}
     \epsilon(\mu,d)\,\sigma_\mu f_{\operatorname{prev}(\mu,d)}.
 \label{eq:ex-trace-probe-observable-average}
\end{align}
This is a real-linear Reynolds average. Reindexing by $\mu\mapsto g\mu$
for each signed motion $g$ and using the defining previous-axis and sign
formulas gives the exact covariance
$\rho_g Z_d=\operatorname{sign}(g,d)Z_{g.\operatorname{perm}(d)}$.
Toggling the second motion coordinate reindexes the conjugation action and
proves the reality condition. In particular, central inversion fixes every
axis and has sign $-1$, so $\rho_{\mathrm{archV}}Z_d=-Z_d$ for each $d$.
Oddness is a consequence of this vector-valued Reynolds projection, not a
claim about arbitrary grade-five labels.

Let $\operatorname{Col}$ be the finite-volume column map from grade-five
labels to coefficient arrays. It is real-linear and commutes with each
motion by finite-volume covariance. The paid labels obey
$\|\operatorname{Col}(X_d)\|_1,\|\operatorname{Col}(Y_d)\|_1\leq2$;
this is the transported column bound, including the norm of the observation
map, rather than a bound on the raw label alone. Define its averaged coefficient cost
by
\[
 c=R\sup_d\sum_{\mu\in\mathfrak A}
       \|\rho_\mu\|_{\mathrm{col}}.
\]
The triangle inequality yields
$\|\operatorname{Col}(\mathscr S_5V)_d\|_1\leq
c\sup_e\|\operatorname{Col}(V_e)\|_1$. This is the precise finite
column estimate used below; the signed or conjugate motion is real-linear,
not assumed complex-linear or isometric in the column norm. The observable average has total
variation one, so, on chamber-interior observables, the reflection seam
makes it preserve a common root-seminorm bound.

Let $x_d,y_d$ be two arrays of finite-volume observables, each interior to
a fixed locality window. Let $X_d,Y_d$ be their grade-five paid
coefficients. The paid receipts and the finite-volume column map give the transported
column bounds $\|\operatorname{Col}(X_d)\|_1\leq2$ and
$\|\operatorname{Col}(Y_d)\|_1\leq2$.
Write $\mathcal O$ for the linear finite-volume observation into a
one-point observable, $\mathcal B$ for pullback from the coarse volume to
the reflection chamber at level $N$, and $\mathcal R$ for the chamber root
seminorm. Assume each paid residual, each pulled-back local field, and each observed
label below is interior to the same reflection chamber, with the locality
window and room conditions required for the finite-volume seam. On this
interior class, assume the seam gives root-seminorm subadditivity and
invariance under each motion. The exact covariance identities are
\begin{equation}
 \mathcal O\,\mathscr S_5=\mathscr S_0\,\mathcal O,
 \qquad
 \mathcal B\,\mathscr S_0=\mathscr S_0\,\mathcal B.
 \label{eq:ex-trace-probe-covariance}
\end{equation}
These follow term by term from linearity and covariance of the finite-volume
observation and pullback under each signed motion.

The paid-component estimates used here are the following explicit
inequalities, for every axis $d$:
\begin{align}
 \mathcal R\!\left(\mathcal B(qx_d)-\mathcal O(a^5X_d)\right)&\leq E,
 &\mathcal R\!\left(\mathcal B(qy_d)-\mathcal O(a^5Y_d)\right)&\leq E.
 \label{eq:ex-trace-probe-paid-errors}
\end{align}
The buffered chamber estimate is this: if
\[
 J_d=q\bigl(x_d-\beta y_d\bigr)
\]
and $g$ is any admissible coarse observable with supremum norm at most
$\tau_j(150000)$, interior to the same locality window, with that window
fitting in the available room, then at any two distinct coarse sites,
\begin{equation}
 \left|\mathbb E\left[(\mathcal B\mathscr S_0J)_d(x)\,
                    (\mathcal B g)(z)\right]\right|
       \leq\tau_j(-30000).
 \label{eq:ex-trace-probe-buffered}
\end{equation}
The estimate is supplied by the finite-volume chamber buffer; its proof is
upstream of this assembly. The scale assumptions and the paid-component
bounds are supplied by the construction's budget and response receipts.

\begin{proposition}[Probe from paid components and a chamber buffer]
\label{prop:ex-trace-probe-assembly}
Under the assumptions above, define
\[
 Z_d=(\mathscr S_5(X-\beta Y))_d,\qquad
 J_d=q(x_d-\beta y_d),\qquad
 P_d=\mathcal O(a^5 Z_d),\qquad
 U_d=(\mathcal B\mathscr S_0J)_d.
\]
Then $Z$ is signed-covariant, with each component satisfying
$\rho_{\mathrm{archV}}Z_d=-Z_d$. Moreover, for every axis $d$,
\begin{enumerate}
\item $\|\operatorname{Col}(a^5 Z_d)\|_1\leq a^5t^4$;
\item $U_d$ is chamber-interior and
  $\mathcal R(U_d-P_d)\leq\tau_j(-200)$;
\item for every admissible bounded local test in
  \eqref{eq:ex-trace-probe-buffered}, the separated-site pairing of $U_d$
  with its pullback is at most $\tau_j(-200)$ in absolute value.
\end{enumerate}
Thus $P_d$ has exactly the coefficient, parity, root-error, and local-price
parts required of the finite-volume probe.
\end{proposition}
\begin{proof}
Real-linearity of the column map and the transported column bounds give
$\|\operatorname{Col}(X_d-\beta Y_d)\|_1\leq2+2\beta$. Covariance
lets the finite-average bound apply to $Z$, so
\[
 \|\operatorname{Col}(Z_d)\|_1
   \leq c(2+2\beta)\leq c(2+4t)\leq6ct\leq t^2.
\]
Since $t\geq1$ we have $t^2\leq t^4$. The column map is real-linear and
the column norm is homogeneous under positive real scalars, so, as
$a^5>0$,
\[
 \|\operatorname{Col}(a^5Z_d)\|_1=a^5\|\operatorname{Col}(Z_d)\|_1
   \leq a^5t^2\leq a^5t^4.
\]
No lower bound on $a$ enters. The finite signed-motion average has the covariance just established;
central inversion fixes each axis with sign $-1$, giving the oddness
assertion for every component.

By linearity and the covariance identities,
\[
 U_d-P_d=\bigl(\mathscr S_0\bigl(\mathcal B(qx)-\mathcal O(a^5X)\bigr)\bigr)_d
       -\beta\bigl(\mathscr S_0\bigl(\mathcal B(qy)-\mathcal O(a^5Y)\bigr)\bigr)_d.
\]
The root seminorm is invariant under each signed motion, is subadditive,
and the averaging weights have total absolute value one. Hence
\[
 \mathcal R(U_d-P_d)\leq(1+\beta)E
   \leq\tau_j(230)\tau_j(-430)=\tau_j(-200),
\]
where $1+\beta\leq1+2t\leq\tau_j(230)$ was used. The interiority of $U_d$
follows because each motion preserves the fixed locality window and finite
averaging preserves the interior domain.

Finally, the buffered estimate gives a separated-site price at most
$\tau_j(-30000)$, which is at most $\tau_j(-200)$ by monotonicity of
$\tau_j$. This proves all three conclusions.
\end{proof}

The proposition isolates the actual finite averaging and error arithmetic
in the selected probe producer. The source's remaining inputs to this
local argument are the two paid-component estimates
\eqref{eq:ex-trace-probe-paid-errors}, the finite-volume covariance
identities, and the buffered chamber estimate
\eqref{eq:ex-trace-probe-buffered}. For all sufficiently large $j$, the
paid-component estimates and the transported column bounds are the
conclusions of Proposition~\ref{prop:ex-trace-bill-root-price} with
grade $5$, applied to each observable and paid label $(x_d,X_d)$ and
$(y_d,Y_d)$; that proposition assumes the paid label bounds and the paid
root receipt of each pair and does not produce them. The buffered estimate is the conclusion of
Proposition~\ref{prop:ex-trace-crown-buffer}, in which $F_j=\tau_j$,
applied to the scale-$k$ field $f=(\mathscr S_0J)_d$. That proposition
requires a Crown witness for this particular averaged combination,
together with its seam, transfer and one-site inputs; the witness is a
separate input of the construction and is produced neither there nor
here. The scale bounds \eqref{eq:ex-trace-probe-scales} are likewise
inputs, holding for all sufficiently large $j$ in the source. In
particular, this section does not claim that arbitrary local components
or an arbitrary chamber record satisfy those paid estimates.
